\section{Introduction}
\label{sec:intro}

SQIsign~\cite{sqisign-spec-v2.0.1} is the only isogeny-based signature scheme
in Round~2 of NIST's additional signature standardization process.
Its 148-byte signatures and 65-byte public keys at NIST-I are smaller
than any other post-quantum scheme under consideration.  But that
compactness comes at a cost: implementing SQIsign requires quaternion
algebras, the Deuring correspondence, dimension-2 theta-coordinate
isogenies, and a signing algorithm that touches all of them.

The specification~\cite{sqisign-spec-v2.0.1} is mathematically precise but
leaves many implementation choices unstated.  The C reference
implementation fills in those choices---but they are embedded in code,
not documented.  An implementor working from the spec alone will make
wrong guesses.  We know because we did.

This paper documents \texttt{sqisign-selkie}, a Rust library for the
NIST-I parameter set ($p = 3 \cdot 2^{324} - 1$) only.  By fixing
the parameter set at compile time, we specialize aggressively: field
arithmetic uses a radix-55 Montgomery representation tuned to this
prime, integer widths match proven worst-case bounds for NIST-I
quaternion intermediates~\cite{EPRINT:KLKL25}, and isogeny degree types are
sized to the 246-bit maximum that NIST-I permits.  We do not attempt
to be generic over parameter sets.

We target four goals:

\begin{enumerate}
  \item \textbf{Interoperability.}  Verify and produce signatures that
    the C reference implementation accepts, byte-for-byte, against the
    NIST KAT vectors.

  \item \textbf{Misuse resistance.}  If a value is always positive and
    odd, it should be impossible to construct one that isn't.  If a
    lattice operation requires Hermite Normal Form, the compiler should
    reject non-HNF input.  This protects not just end-users of the
    public API, but us writing the internals, and anyone who maintains
    this code after us.

  \item \textbf{Constant-time signing (goal, not yet achieved).}
    The specification does not discuss side channels.  The C reference
    is variable-time for the entire quaternion layer.  We traced the
    signing algorithm and confirmed that the quaternion operations
    compute over secret-key-derived data.  Published
    attacks~\cite{EPRINT:MCJKR25} demonstrate practical SPA on
    Cornacchia's algorithm.  Constant-time execution is not a
    nice-to-have; it is a requirement.  Our signing and key generation
    code is currently variable-time, with every secret-touching code
    path marked \texttt{TODO(ct)}.  End-to-end interoperability
    against the C reference is complete (\S\ref{sec:status}); the
    constant-time hardening described in \S\ref{sec:ct} is the
    active workstream.

  \item \textbf{Idiomatic Rust.}  Methods on types over free functions.
    Invariants via distinct types, not runtime checks.  Code that a
    Rust developer would recognize as
    natural~\cite{rust-api-guidelines}.
\end{enumerate}

This paper is a living document, updated as the implementation progresses.

\paragraph{Contributions.}
\begin{itemize}
  \item A catalog of implementation bugs (\S\ref{sec:bugs},
    \S\ref{sec:fixed-width-bugs}) with root causes and lessons.
  \item Compile-time invariant enforcement via Rust's type system
    (\S\ref{sec:types}).
  \item A survey of published constant-time SQIsign techniques and
    the in-progress hardening work
    (\S\ref{sec:ct}, \S\ref{sec:post-baseline}).
  \item A negative test vector suite for SQIsign verification,
    signing, and key generation (\S\ref{sec:testing})---portable
    JSON vectors in C2SP/wycheproof format that other implementations
    can consume to detect vulnerabilities including exponent overflow,
    non-canonical field elements, degenerate kernels, and cross-field
    splicing attacks.
\end{itemize}
